\documentclass[11pt]{article}
\usepackage[T1]{fontenc}
\usepackage{lmodern}
\usepackage[margin=1in,letterpaper]{geometry}
\usepackage{fancyhdr}
\pagestyle{fancy}\fancyhf{}\lhead{Counting zeros from right-hand phases}\rhead{Research note}\cfoot{\thepage}
\setlength{\headheight}{14pt}
\usepackage{amsmath,amssymb,amsthm,booktabs,microtype}
\usepackage[hidelinks]{hyperref}
\newtheorem{theorem}{Theorem}
\newtheorem{lemma}[theorem]{Lemma}
\newtheorem{corollary}[theorem]{Corollary}
\newcommand{\R}{\mathbb R}
\newcommand{\C}{\mathbb C}
\newcommand{\Real}{\operatorname{Re}}
\newcommand{\Imag}{\operatorname{Im}}
\newcommand{\dd}{\,\mathrm d}
\setlength{\emergencystretch}{2em}
\title{Counting critical-line zeros from a fixed strip\\inside the Euler product}
\author{Research note}
\date{26 November 2026}
\begin{document}
\maketitle
\begin{abstract}
Choose any vertical strip of positive width strictly to the right of $\Real s=1$.
For a primitive Dirichlet $L$-function we give a finite test that returns the
exact number of nontrivial zeros in a height interval. When neither endpoint
is a zero ordinate, increasing the number of observation lines eventually
makes the test conclusive. Both the estimate and its error bound use phases
and logarithmic derivatives within the chosen strip. The test needs no
preliminary zero list or lower bound on endpoint distances. The proof combines
the critical-line theorem with a positive rational kernel and a factorization
of its tails. A strip variant certifies one simple zeta zero at heights
between 10 and 18 without using RH. We also give the sharp error-amplification factor for its linear
reconstruction of the unnormalized kernel.
\end{abstract}

\section{Data and the counting theorem}
Let $c=1/2$. For the primitive character of conductor one use
$F(s)=\xi(s)=\tfrac12s(s-1)\pi^{-s/2}\Gamma(s/2)\zeta(s)$. For any other
primitive character $\chi$ of conductor $q$ and parity
$\kappa\in\{0,1\}$, use
\[
 F(s)=(q/\pi)^{(s+\kappa)/2}\Gamma((s+\kappa)/2)L(s,\chi).
\]
These completions are entire of order one. They satisfy
$F(s)=\epsilon\overline{F(1-\bar s)}$, $|\epsilon|=1$. By Dirichlet GRH their
zeros are $c+i\gamma$ with $\gamma\in\R$. Throughout, a sum over $\gamma$
includes multiplicity. For a complex character this convention retains the
actual ordinates of that character with their signs.

Fix $1/2<a<b$, and take any finite set $S\subset[a,b]$ of distinct offsets.
For a nonempty $U\subseteq S$ put
\begin{align}
 g_U(x)&=\prod_{h\in U}\frac{h^2}{x^2+h^2},&
 w_h(U)&=\prod_{k\in U\setminus\{h\}}\frac{k^2}{k^2-h^2},&
 Z_U&=\pi\sum_{h\in U}h w_h(U). \label{eq:weights}
\end{align}
For a continuously unwrapped phase $\theta_h(t)=\arg F(c+h+it)$ write
\[
 P_U(t)=\sum_{h\in U} w_h(U) h\theta'_h(t),\qquad
 C_S(A,B)=\frac1{Z_S}\sum_{h\in S} w_h(S)h
                  \big(\theta_h(B)-\theta_h(A)\big).
\]
An additive constant for each phase cancels. For $A<B$ let $N[A,B]$ count
all ordinates in the closed interval.

\begin{theorem}\label{th:count}
For any partition $S=I\sqcup J$ with both parts nonempty, define
\[
 E_S(A,B)=\frac{Z_J}{2Z_S}\big(P_I(A)+P_I(B)\big).
\]
Then $Z_S>0$, $E_S(A,B)\ge0$, and
\[
             C_S(A,B)-E_S(A,B)\ \le\ N[A,B]\ \le\ C_S(A,B)+E_S(A,B).
\]
A unique integer in this interval gives the exact count.
For any sequence of choices of $m=|S|$ distinct offsets in $[a,b]$, split
them into parts of sizes $\lfloor m/2\rfloor$ and $\lceil m/2\rceil$.
If neither endpoint is a zero ordinate, the error bound tends to zero
exponentially in $m$. In particular the test eventually gives the exact count.
\end{theorem}

\section{Why the signed observations count positive masses}
The passage from one line to a product kernel is an algebraic identity.
The positivity underlying the error bound comes from the location of the zeros.
\begin{lemma}\label{lem:phase}
For all $h>0$,
\[
 h\theta'_h(t)=\sum_\gamma\frac{h^2}{h^2+(t-\gamma)^2}.
\]
The sum converges locally uniformly in $t$. Consequently $P_U(t)=
\sum_\gamma g_U(t-\gamma)$, and $Z_U=\int_{\R}g_U(x)\dd x$.
\end{lemma}
\begin{proof}
For fixed $t$, write $s_t=c+it$. The function
$F(s_t+w)F(s_t-w)$ is even and entire in $w$, so it equals an entire
function $H_t(w^2)$. The order of $H_t$ is at most $1/2$. Its zero multiset
in the variable $z=w^2$ is $\{-(\gamma-t)^2\}_\gamma$.
In particular, a zero of multiplicity $r$ at $s_t$ gives order $2r$ in $w$
and order $r$ in $z$. Equal squared distances add their multiplicities.
Hadamard factorization for an entire function of order less than one therefore gives
\[
 \frac{H'_t(z)}{H_t(z)}
       =\sum_\gamma\frac{1}{z+(\gamma-t)^2}.
\]
The zero count in $|\gamma|\le T$ is $O_q(T\log(2+T))$, which also gives the
asserted local convergence. Evaluate at $z=h^2$ and differentiate the defining product.
The reflection equation gives
$(F'/F)(s_t-h)=-\overline{(F'/F)(s_t+h)}$.
Hence $H'_t(h^2)/H_t(h^2)=\Real(F'/F)(s_t+h)/h=\theta'_h(t)/h$.
Finally,
\[
 g_U(x)=\sum_{h\in U} w_h(U)\frac{h^2}{x^2+h^2}.
\]
The residues agree as rational functions of $x^2$, and both sides vanish
at infinity. Integration evaluates $Z_U$.
\end{proof}

\begin{proof}[Proof of Theorem~\ref{th:count}]
For $d\ge0$ put $Q_S(d)=Z_S^{-1}\int_d^\infty g_S(u)\dd u$.
All factors decrease on $[0,\infty)$, so for $d,y\ge0$,
\[
 g_S(d+y)\le g_I(d)g_J(y),\qquad
 Q_S(d)\le \frac{Z_J}{2Z_S}g_I(d).                 \tag{1}
\]
Consider a single ordinate $\gamma$. Its contribution to $C_S$ is the mass
of $[A,B]$ for the probability density $g_S(t-\gamma)/Z_S$.
If $A\le\gamma\le B$, the missing mass is the sum of the two exterior
tails. If $\gamma$ lies outside, the included mass is at most the tail at
the nearer endpoint. In both cases the absolute error is bounded by
$Q_S(|A-\gamma|)+Q_S(|B-\gamma|)$. This bound includes equality of an
ordinate and an endpoint, for which $Q_S(0)=1/2$. Sum (1) over ordinates
and apply Lemma~\ref{lem:phase}. The exchange with the interval integral is
justified by nonnegativity.

For exponential convergence write
$D_k=\int_\R(1+u^2)^{-k}\dd u
=\sqrt\pi\,\Gamma(k-1/2)/\Gamma(k)$. Direct comparison gives
$aD_{|U|}\le Z_U\le bD_{|U|}$. Since $D_k\sim\sqrt{\pi/k}$,
$Z_J/Z_S$ is bounded for the prescribed split. At an endpoint $e$ at distance
$d_e>0$ from the zero set,
\[
  P_I(e)\le (1+d_e^2/b^2)^{-(|I|-1)}
       \sum_\gamma \frac{b^2}{b^2+(e-\gamma)^2}.
\]
The last sum is finite. This proves the rate and termination. The computed
bound $E_S$ itself requires no value of $d_e$.
\end{proof}

\section{Evaluation and its cost}
When $c+h>1$, a continuous logarithm of the $L$-factor is available from
\[
  \log L(s,\chi)=\sum_p\sum_{k\ge1}\frac{\chi(p)^k}{k p^{ks}},
  \qquad \frac{L'}{L}(s,\chi)=-\sum_{n\ge1}\frac{\chi(n)\Lambda(n)}{n^s}.
\]
Both series converge absolutely. The gamma logarithm on $\Real z>0$ is
normalized on the positive real axis. For $\xi$, the elementary factors
$s$ and $s-1$ also have continuous logarithms in the observation region.
Thus $\theta_h(B)-\theta_h(A)$ can be obtained by subtraction of these
specified logarithms at the two endpoints. Taking a principal argument
of the ratio of the endpoint values can lose whole turns. With numerical
data one encloses $C_S$ and $E_S$ before checking uniqueness of the integer.

\paragraph{A count from ten endpoint evaluations.}
For $F=\xi$ take $[A,B]=[10,18]$, $I=\{5/2,7/2,9/2\}$ and $J=\{3,4\}$.
Thus every observed real part is between $3$ and $5$. The enclosed values
from \texttt{phase\_check.py}, rounded outward here, are
\[
 1.00585<C_S<1.00634,\qquad 0.28792<E_S<0.28922 .
\]
The resulting enclosure lies inside $(0.7166,1.2956)$, whose only integer
is $1$. This example uses prime powers through $256$ in the two absolutely
convergent series. For real part $\sigma$ the respective tails are bounded by
$D=256^{1-\sigma}/(\sigma-1)$ and
$D(\log 256+1/(\sigma-1))$. The script also bounds the Stirling remainders
for the gamma phase and slope. All numerical inputs to the count are computed
without supplying any zero ordinate.

The coefficients can be large even when the final count is small. We quantify
that issue for the reconstruction of $P_S$ from the scaled slopes
$y_h=h\theta'_h(t)$. A common absolute error $\eta$ in these inputs gives a
worst-case error $\eta\sum_h |w_h(S)|$.
\begin{theorem}\label{th:minimax}
Among all $m\ge2$ distinct offsets in $[a,b]$, the minimum amplification is
\[
 T_{m-1}\!\left(\frac{b^2+a^2}{b^2-a^2}\right)
 =\frac12\left[\left(\frac{b+a}{b-a}\right)^{m-1}
               +\left(\frac{b-a}{b+a}\right)^{m-1}\right],
\]
where $T_n(\cos\vartheta)=\cos(n\vartheta)$. It is attained by taking
$h^2$ to be the images of the $m$ extrema of $T_{m-1}$ in $[a^2,b^2]$.
\end{theorem}
\begin{proof}
The $w_h$ are Lagrange interpolation weights for evaluation at $0$ from
nodes $h^2$. Apply the formula to $T_{m-1}$ composed with the affine map from
$[a^2,b^2]$ onto $[-1,1]$. Its modulus on the nodes is at most one, giving
the lower bound by the triangle inequality. At its ordered extrema the
polynomial alternates in sign, and so do the interpolation weights since
$0$ lies strictly to the left of the nodes. All terms of the interpolation
sum then have the same sign. The displayed expression follows from the
formula for $T_n$ outside $[-1,1]$.
\end{proof}
This minimization treats independent absolute errors on the scaled slopes.
Phase differences, gamma factors and normalization may have different
conditioning in an implementation.

\section{A finite test with no critical-line assumption}
The same data can sometimes certify a count using only the classical strip
$|\beta-c|\le D=1/2$ for each zero $\beta+i\gamma$. Write $N[A,B]$ for
the number of zeros with imaginary part in the interval, at any real part.
\begin{theorem}\label{th:strip}
Define
\[
 K_S=\prod_{h\in S}(1-D^2/h^2)^{-1},\quad
 \alpha_I=\sum_{h\in I}\arcsin(D/h),\quad
 L_I=\cos(\alpha_I)\prod_{h\in I}(1+D^2/h^2)^{-1}.
\]
If $\alpha_I<\pi/2$, then $P_I(t)\ge0$ and
\[
             |N[A,B]-C_S(A,B)|\le (K_S/L_I)E_S(A,B).
\]
\end{theorem}
\begin{proof}
Apply the same order-$1/2$ product to $H_t$ without fixing the real parts of
its zeros. For each $\rho=\beta+i\gamma$ put $\delta=\beta-c$.
Since $(\rho-c-it)^2=-(t-\gamma+i\delta)^2$, the partial-fraction identity gives
$P_U(t)=\sum_\rho\Real g_U(t-\gamma+i\delta)$. For real $x$,
\[
 |g_S(x+i\delta)|\le K_S g_S(x),\qquad
 \Real g_I(x+i\delta)\ge L_I g_I(x).
\]
Indeed $|h^2+(x+i\delta)^2|$ is bounded below by
$(1-D^2/h^2)(h^2+x^2)$ using the real part, and bounded above by
$h^2+x^2+D^2$ by the triangle inequality. Its absolute argument is at most
$\arcsin(D/h)$, since $2|x\delta|/(h^2-\delta^2+x^2)$ is at most
$|\delta|/\sqrt{h^2-\delta^2}$. Summing these angles proves the second
inequality. The integral of $g_S(x+i\delta)$ over the whole horizontal
line equals $Z_S$ by a contour shift, since its poles have imaginary parts
$\pm h$ with $h>D$. The omitted or included tails for a single zero are
therefore bounded in absolute value by $K_S$ times the two real tails in
the proof of Theorem~\ref{th:count}. Summing and using
$\sum_\rho g_I(e-\gamma)\le P_I(e)/L_I$ proves the result.
\end{proof}
For the five lines above, $K_S/L_I<1.347$. The count enclosure still lies
strictly inside $(0,2)$, so exactly one zeta zero has height in $[10,18]$
even without invoking RH. In fact the zero is simple and lies on the
critical line: reflection sends any off-line zero to a distinct one at the
same height. For intervals with arbitrarily narrow endpoint clearances,
the eventual-success statement of Theorem~\ref{th:count} uses GRH; the angle
condition in Theorem~\ref{th:strip} restricts the allowable size of $I$.

\section{What a single line records before assuming GRH}
The squared variable in Lemma~\ref{lem:phase} also gives an unconditional
identity. Use the same $H_t$ and let $r_h(t)$ be the radius of convergence
of the Taylor series of $H'_t/H_t$ about $z=h^2$, with $h>1/2$. All of its
Taylor coefficients can be obtained by real differentiation of
$\Real(F'/F)(c+\sqrt z+it)/\sqrt z$ at that positive point.
\begin{corollary}
Writing $\rho=\beta+i\gamma$ for the nontrivial zeros and
$\Delta=\sup_\rho |\beta-c|$, one has
\[
 \inf_{t\in\R} r_h(t)=h^2-\Delta^2.
\]
Under Dirichlet GRH the stronger pointwise formula is
$r_h(t)=h^2+\min_\gamma(t-\gamma)^2$.
\end{corollary}
\begin{proof}
The poles of $H'_t/H_t$ are precisely $z=(\rho-c-it)^2$. Coincident poles
have positive integral residues and cannot cancel. The Taylor radius is the
distance to the nearest pole. For $u=\beta-c$, $v=\gamma-t$,
\[
 |h^2-(u+iv)^2|^2=(h^2-u^2)^2+2(h^2+u^2)v^2+v^4.
\]
For each zero this is minimized over $t$ at $t=\gamma$. Since
$|u|\le1/2<h$, taking infima gives the first formula. Taking $u=0$ for
every zero gives the second.
\end{proof}

\section*{Sources and relation to earlier work}
The formal input is the critical-strip zero theorem in
\texttt{grh/src/Main.lean} of the Sernau repository~\cite{lean}.
The linked overview~\cite{wiki} explains the explicit-formula connection
between prime oscillations and zeros. The proof above uses the output of
the formal development and the classical functional equation and Hadamard
product. In particular the reflection step works for complex primitive
characters individually. The proofs in this note are given in ordinary
mathematical form.

Single-shift phase functions for completed $L$-functions occur in
Lagarias~\cite{lagarias}. Rational combinations of resolvents for smoothing
spectral measures are developed by Colbrook, Horning and Townsend~\cite{cht}.
B\"uthe~\cite{buthe} uses explicit-formula test functions to verify completeness
of zero lists. The extremal polynomial argument belongs to classical
Chebyshev extrapolation~\cite{hl}. The specific contribution proposed here
is the computable two-endpoint bound (1) combined with fixed-strip termination
and the conditioning calculation for these observations. These comparisons
establish the context for that claim. Bibliographic priority remains open.

\begin{thebibliography}{9}
\bibitem{lean} L. Sernau et al., \emph{GRH for Dirichlet $L$-functions, in Lean 4},
commit \texttt{72fead10}, 2026.
\href{https://github.com/lsernau/grh-lean/tree/72fead102f663d5c7a407964a13dc0cf9d977207}{Pinned repository}.
A. Sutherland, \href{https://math.mit.edu/~drew/shortgrh/}{independent verification record}.
\bibitem{wiki} Wikipedia contributors, \emph{Riemann hypothesis}, consulted
26 November 2026, \url{https://en.wikipedia.org/wiki/Riemann_hypothesis}.
\bibitem{lagarias} J. C. Lagarias, \emph{Zero spacing distributions for
differenced $L$-functions}, \href{https://arxiv.org/abs/math/0601653}{arXiv:math/0601653}.
\bibitem{cht} M. J. Colbrook, A. Horning and A. Townsend,
\emph{Computing spectral measures of self-adjoint operators},
SIAM Rev. 63 (2021), 489--524, \href{https://arxiv.org/abs/2006.01766}{arXiv:2006.01766}.
\bibitem{buthe} J. B\"uthe, \emph{A method for proving the completeness of
a list of zeros of certain $L$-functions}, \href{https://arxiv.org/abs/1308.6704}{arXiv:1308.6704}.
\bibitem{hl} P. G. Hoel and A. Levine, \emph{Optimal spacing and weighting
in polynomial prediction}, Ann. Math. Statist. 35 (1964), 1553--1560.
\end{thebibliography}
\end{document}
